\documentclass[twocolumn]{article}
\usepackage{graphicx}
\usepackage{amsmath}
\usepackage{hyperref}
\usepackage[utf8]{inputenc}
\usepackage{geometry}
\geometry{a4paper, margin=1in}

\providecommand{\tightlist}{%
  \setlength{\itemsep}{0pt}\setlength{\parskip}{0pt}}

$if(highlighting-macros)$
$highlighting-macros$
$endif$

\title{Autonomous Inverse Design of Helical Inertial Filters for Lunar Exploration: A Software-Defined Engineering Approach}
\author{Lawrence Kincheloe \\ \textit{NASA / Open Source Community}}
\date{\today}

\begin{document}

\maketitle

\begin{abstract}
As humanity prepares for sustained lunar presence, the mitigation of abrasive, electrostatic lunar regolith poses a critical challenge to Environmental Control and Life Support Systems (ECLSS). Traditional barrier filtration methods are prone to rapid clogging in such environments. This paper presents the OpenAuto-CFD "Software-Defined Engineering" framework for the autonomous design and optimization of fluid dynamics components, and applies it to helical inertial separators ("corkscrew filters"). By integrating parametric Computer-Aided Design (OpenSCAD), Computational Fluid Dynamics (OpenFOAM), and Generative Artificial Intelligence (LLM), the OpenAuto-CFD system performs an automated search of the high-dimensional design space. We demonstrate that this inverse design methodology can effectively navigate the complex fluid dynamic regime defined by Reynolds ($Re$), Dean ($De$), and Stokes ($Stk$) numbers. Experimental validation confirms the theoretical predictions, highlighting the critical dependence of separation efficiency on inlet velocity and the generation of stable Dean vortices.
\end{abstract}

$body$

\end{document}
